\chapter*{Manual Outline}
\addcontentsline{toc}{chapter}{Manual Outline}

This manual is organized as a chaptered technical document. The goal is to make
the project explainable, reproducible, and easy to extend chapter by chapter.

\begin{table}[htbp]
\centering
\caption{Planned chapter structure.}
\begin{tabular}{@{}>{\raggedright\arraybackslash}p{0.08\linewidth}
>{\raggedright\arraybackslash}p{0.26\linewidth}
>{\raggedright\arraybackslash}p{0.54\linewidth}@{}}
\toprule
No. & Chapter & Scope \\
\midrule
1 & Introduction & Research motivation, workspace goals, intended users, and the
problem TEDEX solves. \\
2 & Workspace Architecture & Folder layout, source-code boundaries, data
ownership, and generated-output conventions. \\
3 & Raw TE Analysis & ZEM/LFA ingestion, metadata requirements, transport
property calculation, and extracted feature files. \\
4 & SPB Fitting & Effective-mass fitting, performance curves, conductivity-axis
fitting, command templates, outputs, and interpretation. \\
5 & Plotting and Figures & Standard TE figures, flexible plotting recipes,
style conventions, and publication-ready export formats. \\
6 & XRD Analysis & XRD data layout, normalized/stacked plots, PDF-card overlays,
and lattice-parameter fitting. \\
7 & Assessment and Prediction & Batch assessment, AI-assisted interpretation,
Bayesian prediction, and follow-up candidate ranking. \\
8 & Reproducible Operations & End-to-end batch protocol, logging, figure slots,
and recommended archival practice. \\
\bottomrule
\end{tabular}
\end{table}

\section*{Editing Convention}

Each chapter should include three practical components:
\begin{enumerate}[leftmargin=*]
  \item a conceptual description of the workflow;
  \item command slots with exact bash commands;
  \item terminal-output and figure slots that can be filled after each real run.
\end{enumerate}
